% !TEX root = main.tex

% ========================================================
% BIBLIOGRAPHIE
% ========================================================
\newpage
\addcontentsline{toc}{chapter}{Bibliographie et Nétographie}
\begin{thebibliography}{99}

\bibitem{machining} 
\textbf{Smid, Peter.} \textit{CNC Programming Handbook, Third Edition}. Industrial Press Inc., 2007. \newline
(Référence canonique sur la structure et le parsing du G-Code ISO 6983).

\bibitem{robotics} 
\textbf{Craig, John J.} \textit{Introduction to Robotics: Mechanics and Control}. Pearson Prentice Hall, 2005. \newline
(Méthode de Denavit-Hartenberg et calculs matriciels de cinématique inverse).

\bibitem{rdm} 
\textbf{Timoshenko, Stephen P.} \textit{Résistance des matériaux}. Dunod. \newline
(Calcul des flèches sur les guides de translation HGR15).

\bibitem{esp32} 
\textbf{Espressif Systems.} \textit{ESP32 Technical Reference Manual}, 2023. \newline
\texttt{https://www.espressif.com/en/support/documents/technical-documents} \newline
(Architecture Dual-Core et FreeRTOS).

\bibitem{flutter} 
\textbf{Dev Google.} \textit{Flutter Architectural Overview}. \newline
\texttt{https://docs.flutter.dev/resources/architectural-overview} \newline
(Documentation sur le Render Tree, StateNotifier et le pattern Riverpod).

\bibitem{openbuilds}
\textbf{OpenBuilds Community.} \textit{Open Source Machine Design Archives}. \newline
\texttt{https://openbuilds.com/} \newline
(Inspiration pour les sous-assemblages de la géométrie V-Slot des profils aluminium).

\end{thebibliography}

% ========================================================
% ANNEXES
% ========================================================
\newpage
\appendix
\chapter{Annexes techniques}

\section{Extrait de l'algorithme "Look-Ahead" du Planner (C++)}
Cet annexe présente une portion simplifiée de l'accélération trapézoïdale programmée sur le Core 1 de l'ESP32 pour lisser les micro-segments du G-Code 5 axes.

\begin{tcolorbox}[colback=white, colframe=macouleur, title=Planner.cpp (Extrait)]
\begin{verbatim}
void plan_buffer_line(float x, float y, float z, float a, float b, float feed_rate) {
    // Calcul du vecteur déplacement 5D (Espace Machine)
    float dx = x - position[X_AXIS];
    float dy = y - position[Y_AXIS];
    float dz = z - position[Z_AXIS];
    float da = a - position[A_AXIS];
    float db = b - position[B_AXIS];

    float distance = sqrt(dx*dx + dy*dy + dz*dz + da*da + db*db);

    // Blocage si déplacement inférieur à la tolérance mécanique (0.01mm)
    if (distance < MINIMUM_PLANNER_SPEED) return;

    // Calcul de la vitesse de croisière cible (Target Speed)
    float nominal_speed = feed_rate;
    
    // Détection d'angle aigu (Jonction) entre l'ancien bloc et le nouveau
    float junction_cos_theta = calculate_junction_deviation(previous_vector,
                                                            current_vector);
                                                            
    // Ralentissement dynamique aux angles (Cornering Speed)
    float max_junction_speed = min(nominal_speed, 
                                   sqrt(JUNCTION_DEVIATION * acceleration 
                                   / (1.0 - junction_cos_theta)));

    // Insertion dans le Ring Buffer
    block_buffer[head].nominal_speed = nominal_speed;
    block_buffer[head].entry_speed = max_junction_speed;
    
    // Avance de l'index du Buffer Circulaire
    head = (head + 1) % BLOCK_BUFFER_SIZE;
}
\end{verbatim}
\end{tcolorbox}
